\chapter{信息渠道与工具}
\label{ch21}
\ist{信息与工具}
\begin{keybox}
本章处理推免的共同前提：你从哪儿知道、文件怎么读、消息怎么辨真假。
\textbf{信息渠道}回答去哪个官网的哪个栏目、国家平台与学校官网冲突时以谁为准、往年拟录取名单在哪里找；\textbf{信息管理}给出台账与跟踪表的字段设计、归档命名规则和防漏看邮件的做法；\textbf{文件解读}专门处理「原则上」「同等条件下优先」「以实际下达为准」这类表述的读法与例外条件；\textbf{人脉渠道}讲怎么找学长学姐、怎么问才拿得到真话，以及给导师发邮件的组织方式；\textbf{数据分析}教你怎么用招生专业目录、拟录取公示与历年数据自己算名额与竞争烈度；\textbf{信息甄别}给出群消息、公众号、经验帖的可信度排序与核实动作；\textbf{导师调研}只看公开可核验的信号；\textbf{系统操作}讲推免服务系统的确认时限与常见误操作；\textbf{机构服务}守住底线——不存在能操控录取结果的合法服务，凡称「内部名额」「保过通道」「包录取」的均不可信。
阅读方式：全国统一规定写条文与文号；院校规定点明学校与年份，不代表你校口径；查不到的如实标「未见规定」，并告诉你去问谁。
\end{keybox}

\section{信息管理}
\begin{qa}{S0387}{个人保研信息台账的表格应该包含哪些字段？}
台账的作用是让「我不知道自己卡在哪一步」这件事消失。建议建三张表，字段直接对齐官方流程。
\begin{itemize}
  \item 院校目标表：学校、学院、专业/方向、招生类型（硕士/直博）、拟招生人数及出处、报名系统入口、是否需预报名、往年公示链接、咨询电话与邮箱。拟招生人数在招生专业目录里逐专业列出，如华北电力大学2027年推荐免试研究生招生专业目录按学院、联系方式、招生类型、专业代码、学位类型、专业名称、拟招生人数、学习方式、研究方向展开 \nref{S0414}。
  \item 投递跟踪表：志愿序号、提交时间、可修改时间（系统规则是每个志愿提交后48小时内不允许修改）\nref{S0444}、当前状态（未提交/待复试/已拒/待录取）、回复截止时间、该状态依据哪份文件。
  \item 材料与节点表：材料名称、版本、盖章要求、提交渠道与截止时间、快递单号。东北师范大学要求把多份材料合并为一个 PDF 上传至「佐证材料」栏并按规定命名 \nref{S0607}，此类格式要求必须写进表里。
\end{itemize}
\nref{S0414}
\nref{S0444}
\nref{S0607}
\end{qa}

\iin{华北电力大学}
\ith{信息管理}
\ist{信息与工具}
\begin{qa}{S0388}{保研全流程有哪些容易错过的时间节点，怎么做提醒？}
系统硬节点是公开的，容易被漏的都是学校自设的软节点。以2027年为例：9月18日起注册填报个人资料并查看本人推免资格；9月21日9:00开始填报志愿；10月18日截止填报志愿；9月22日8:00开始确认招生单位发送的复试通知；9月24日14:00开始确认待录取通知；10月20日关闭系统 \nref{S0464}。学校层面的开关时间更细，河北大学公布的系统功能时间表为：9月21日9:00开通报名、9月22日8:00开通复试通知与确认、9月24日8:00发送待录取、9月24日14:00开通待录取确认、10月18日关闭注册支付报名、10月20日关闭确认与发送功能 \nref{S0604}。
提醒做法：每个节点写成「日期＋时点＋我要做什么」三列，提前3天和1天各提醒一次；学院预报名与材料寄送截止通常早于系统节点两周以上，必须单列提醒。
\nref{S0464}
\nref{S0604}
\end{qa}

\iin{教育部学生服务与素质发展中心（学信网/研招网）}
\ith{信息管理}
\ist{信息与工具}
\begin{qa}{S0389}{需不需要专门做院校跟踪表来管理投递进度？}
需要。不是因为显得专业，而是系统规则本身会制造状态混乱。
系统规则导致的复杂状态至少有这些：可同时填报三个平行志愿（不分主次），每个志愿提交后48小时内不允许修改，期间招生单位下载报名信息并决定是否允许参加复试；志愿提交48小时后仍未接到复试通知、复试未通过或拒绝待录取的，可继续填报其他志愿；48小时内如招生单位通过系统明确拒绝，可立即填报其他志愿 \nref{S0444}。待录取阶段更严格：待录取通知可拒绝多个但只能接受一个，一旦接受就不能再填报其他志愿，也不能接受其他复试与待录取通知，其他单位也不能再向你发放 \nref{S0444}。
再叠加学校自设的独立报名系统（中国矿业大学有研究生推免预报名系统 \nref{S0626}、南京航空航天大学有研究生报考服务系统 \nref{S0648}），同一所学校可能同时存在两条线，不记表必然记错。
\nref{S0444}
\nref{S0626}
\nref{S0648}
\end{qa}

\iin{教育部学生服务与素质发展中心（学信网/研招网）}
\ith{信息管理}
\ist{信息与工具}
\begin{qa}{S0390}{怎么系统收集和归档各校招生简章，避免往年信息混淆？}
归档的核心是「文件名自带年份和类型」，并且每年新建目录。
\begin{itemize}
  \item 目录结构：按「年度／学校／文件类型」三层。文件类型固定为六类：招生简章、招生专业目录、推免接收办法、学院实施细则、拟录取公示、补充通知。
  \item 命名规则：年份-学校-学院-文件类型-发布日，例如「2027-XX大学-机械学院-接收推免办法-20260908」。发布日在页面上都有，如上海交通大学机械与动力工程学院预推免办法发布时间2026-09-08 \nref{S0262}、中国矿业大学接收通知时间2026-09-08 \nref{S0626}。
  \item 每份文件必须同时存：原始网址、页面截图、下载时间。公示类文件过公示期可能下线，南京理工大学的拟录取名单就是以附件形式发布的 \nref{S0181}。
\end{itemize}
归档对象清单来自制度要求：教学厅〔2013〕8号要求推荐单位与接收单位在本单位网站公布推免名额与推免办法并统一公示 \nref{S0047}；教学厅〔2019〕2号要求公开招生章程、招生政策规定、招生专业目录、分专业招生计划、复试录取办法 \nref{S0431}。
\nref{S0262}
\nref{S0626}
\nref{S0181}
\nref{S0047}
\nref{S0431}
\end{qa}

\iin{上海交通大学}
\ith{信息管理}
\ist{信息与工具}
\begin{qa}{S0403}{保研信息搜集应该从什么时候开始，怎么分阶段做？}
按三段切，越早越省力。
\begin{itemize}
  \item 大一至大二：搞清规则。把本校推免办法读完，确认四件事——名额怎么分、排名课程有哪些、加分认定区间、有无数理与英语等硬门槛。西南交通大学明确各学院须确定专业免研排名课程，并于推免工作当年6月前通过学院官网向学生公布 \nref{S0738}；学校推免指标分配情况通过教务网公布，各学院再通过学院官网公布本专业推免指标分配情况 \nref{S0738}。
  \item 大二暑假至大三上：锁定目标与导师。各校接收信息集中在6—9月分批发布，如华东师范大学2026年7月24日发布预报名通知、9月15日发布综合考核细则汇总 \nref{S0611}\nref{S0613}。
  \item 大三下至9月：执行。系统节点集中在9月中下旬至10月中下旬：9月21日9:00开始填报志愿、10月18日截止填报志愿、10月20日关闭系统 \nref{S0464}。
\end{itemize}
\nref{S0738}
\nref{S0611}
\nref{S0613}
\nref{S0464}
\end{qa}

\iin{西南交通大学}
\ith{信息管理}
\ist{信息与工具}
\begin{qa}{S0404}{大二和大三学生做保研信息准备的重点有什么不同？}
一句话：大二对内（规则与成绩结构），大三对外（目标与导师）。
\begin{itemize}
  \item 大二重点：读懂本校推免办法与学院细则；确认排名课程清单，西南交通大学要求各学院确定专业免研排名课程并于当年6月前通过学院官网公布 \nref{S0738}；弄清加分认定的硬边界——多数学校限定学术专长原则上仅限本科阶段在核心期刊上以独立作者或第一作者发表的与学业相关论文，或作为主力成员参加与学业相关的国内权威科研竞赛（全国赛）并获得三等奖以上奖励，且与直系亲属或学历职称明显高于本人者合作的成果不计入 \nref{S0527}\nref{S0738}。
  \item 大三重点：目标调研与投递。夏季起盯夏令营与预推免（各校6—9月分批发布 \nref{S0611}\nref{S0613}）；9月起盯系统（9月21日填报志愿、10月18日截止填报、10月20日关闭）\nref{S0464}；同时按学校清单准备材料格式（东北师范大学要求多份材料合并为一个 PDF 上传）\nref{S0607}。
\end{itemize}
\nref{S0738}
\nref{S0527}
\nref{S0611}
\nref{S0613}
\nref{S0464}
\nref{S0607}
\end{qa}

\iin{西南交通大学}
\ith{信息管理}
\ist{信息与工具}
\begin{qa}{S0408}{加分材料提交的时间节点和常见退回原因有哪些？}
时间节点不在9月。加分认定必须在推免名单产生之前完成，由学院在本校推免工作通知规定的时点集中受理：西南交通大学明确各学院对学生提交的加分申请进行综合评议并确定加分，通过学院官网向学生公布 \nref{S0738}；南京工业大学的做法是推免生初选名单确定后采取多种途径公示，公示内容包括姓名、推免成绩排名（含学业综合成绩与全面发展考核情况）、推荐顺序 \nref{S0527}，可见加分必须赶在初选之前。
常见退回原因（均有制度对应）：
\begin{itemize}
  \item 成果与学业无关、非核心期刊、非独立/第一作者：学术专长原则上仅限核心期刊上以独立作者或第一作者发表的与学业相关论文，或国内权威科研竞赛全国赛三等奖以上 \nref{S0527}\nref{S0738}；
  \item 与直系亲属或学历、职称、职务明显高于本人者合作：仅作为参考，不计入综合评议加分 \nref{S0738}；
  \item 未通过专家审核小组鉴定或未参加公开答辩：未通过的不纳入综合成绩计算体系 \nref{S0527}；
  \item 形式问题：需原件扫描、按清单顺序命名并合并为指定格式（东北师范大学要求多份材料合并为一个 PDF 上传至「佐证材料」栏）\nref{S0607}。
\end{itemize}
\nref{S0738}
\nref{S0527}
\nref{S0607}
\end{qa}

\iin{西南交通大学}
\ith{信息管理}
\ist{信息与工具}
\begin{qa}{S0409}{保研排名公示期间应该核对哪些信息？}
公示期是唯一有效的纠错窗口，要核的是会算进结果的每一项。
国家要求：推荐单位和接收单位都要在本单位网站开辟专栏，公布推免生名额与推免办法，并对拟推免名单、拟接收名单统一公示，公示时间不少于10个工作日；公示期间公示内容不得修改，如有变动，须对变动部分做出说明，并对变动内容另行公示10个工作日 \nref{S0047}\nref{S0318}。公示内容的最低要求见南京工业大学：应包括学生姓名、推免成绩排名（含学业综合成绩与全面发展考核情况）、推荐顺序，公示期间内容不得随意变动 \nref{S0527}。
可核对字段的实例是湖南大学计算机学院2027届公示：学院、专业代码、专业名称、姓名、学号、性别、学业成绩、所在专业学业成绩排名、综合成绩、所在专业综合成绩排名、排名、总人数、GPA、备注（如「基础保研」）\nref{S0131}。据此逐项核：GPA 与成绩单是否一致、排名总人数是否是今年的、备注是否漏标你的专项类型。
同时留证据：公示页面不保证长期保留，名单多以附件发布 \nref{S0181}，公示期内务必下载。
特别提醒「国家规定」与「院校实际」要分开看：国家统一要求是公示时间不少于10个工作日 \nref{S0047}，但实际执行有差异——中国矿业大学写明拟录取的推免生名单在研究生招生网公示不少于7日 \nref{S0626}，湖南大学计算机学院2027届推免名单公示写明「公示期5天」\nref{S0131}。所以不要用10个工作日去推算自己的纠错窗口，要在本校通知里找实际天数。
\nref{S0047}
\nref{S0318}
\nref{S0527}
\nref{S0131}
\nref{S0181}
\nref{S0626}
\end{qa}

\iin{中华人民共和国教育部}
\ith{信息管理}
\ist{信息与工具}
\begin{qa}{S0410}{对排名或加分有异议时的申诉渠道和流程是什么？}
分三层，按顺序走，不要跳级。
第一层，向推荐高校（本校）内部提出。教学厅〔2020〕12号要求将学生对推免程序、标准的意见和对结果的异议纳入本单位信访举报、申诉受理机制，充分保证学生反映问题以及申诉渠道畅通 \nref{S0318}。成都大学商学院的落地做法可作模板：学院设推免工作举报电话和邮箱并向全校公布，学生对遴选推荐工作有异议可向所在学院反映，学院无法解决的也可直接向教务处反映，涉嫌违纪违规的由纪检监察部门依纪依规调查，受理部门应及时查明情况并公布处理结果 \nref{S0079}。
第二层，向招生单位（接收方）提出。考生认为所报考招生单位的招生录取行为违反规定的，可向报考招生单位提出异议、申诉或举报，招生单位应当进行调查处理；属于对政策执行存在异议的，应当及时书面或口头答复申诉人 \nref{S0442}。中国矿业大学公开了研招办与监察处两套电话邮箱作为咨询、监督与申诉渠道 \nref{S0626}；重庆大学数统学院规定成绩公示3日内接受考生书面实名申诉 \nref{S0618}。
第三层，对招生单位书面答复不服的，可向招生单位所在地省级教育行政部门或省级教育招生考试机构申请复查；对复查结论不服的，可按相关规定向上一级机关提出复核 \nref{S0442}。
\nref{S0318}
\nref{S0079}
\nref{S0442}
\nref{S0626}
\nref{S0618}
\end{qa}

\iin{shpg（高校官网）}
\ith{信息管理}
\ist{信息与工具}
\begin{qa}{S0414}{面试通知、材料补交邮件容易进垃圾箱，怎么避免漏看？}
学校文件里已写明「漏看后果自负」，所以这是你的责任。东北师范大学写明考生应及时关注我校及招生院部发布、发送的相关文件与通知，因考生个人疏忽等原因造成的一切后果由考生本人承担 \nref{S0607}；青岛大学提示请考生务必保持通讯畅通 \nref{S0693}；华东师范大学心理与认知科学学院写明具体安排以邮件通知为准，请推免生注意查收邮件、及时回复是否参加，否则视为放弃 \nref{S0295}。
防漏做法四条：一是把学校域名加白名单，研招办邮箱多为 yzb@、yjzs@ 形式（如东北师大 yzb@nenu.edu.cn \nref{S0607}、中山大学 yzba@mail.sysu.edu.cn \nref{S0166}）；二是留两个邮箱并全部开启转发；三是走官方微信服务号与系统内通知双通道（如「东师研招」\nref{S0607}）；四是手机号保持可用，因为学校会通过推免服务系统、电话或电子邮件发送复试相关信息 \nref{S0067}。
\nref{S0607}
\nref{S0693}
\nref{S0295}
\nref{S0067}
\nref{S0166}
\end{qa}

\iin{东北师范大学}
\ith{信息管理}
\ist{信息与工具}
\begin{qa}{S0421}{保研信息过载时，怎么筛选出与自己相关的关键信息？}
用一张四问过滤表替代「都看一遍」。
\begin{itemize}
  \item 问一：它管的是推荐还是接收？推荐类看本校教务处，接收类看目标校研究生院与学院；两类信息的责任主体不同 \nref{S0047}。
  \item 问二：它说的是哪一年、哪个专业方向、哪种类型（硕士／直博／专项）？招生专业目录按专业逐个列出拟招生人数与研究方向（华北电力大学2027年目录列到研究方向层级）\nref{S0414}，不看专业直接看数字必然误判。
  \item 问三：它是硬时限还是软提示？硬时限是系统与学校的日期（9月21日9:00填报志愿、10月18日截止填报、10月20日关闭系统）\nref{S0464}，软提示是宣讲与建议；只对硬时限设提醒。
  \item 问四：它有没有给出一手文件出处？给不出的一律进「待核」清单。
\end{itemize}
人群分层：边缘生优先盯公示与申诉机制（公示期不少于10个工作日、变动须另行公示）\nref{S0318}；双非学生优先核对接收方公开的申请条件，部分学院把「毕业高校为『双一流』建设高校及建设学科或国家重点学科」写进条件 \nref{S0262}。
\nref{S0047}
\nref{S0414}
\nref{S0464}
\nref{S0318}
\nref{S0262}
\end{qa}

\iin{中华人民共和国教育部}
\ith{信息管理}
\ist{信息与工具}
\section{人脉渠道}
\begin{qa}{S0379}{怎么找到目标院校目标学院的学长学姐？}
官方渠道其实会主动创造接触机会，先走官方再走私人。
\begin{itemize}
  \item 夏令营与开放日：学院常把「与学长学姐面对面」写进活动内容。天津大学软件学院在夏令营活动内容里明确「组织实验室交流活动，为营员搭建与导师、学长学姐面对面互动的平台」\nref{S0306}；天津大学化工学院夏令营也安排学生座谈与学术汇报交流 \nref{S0304}。这是最正当的接触窗口。
  \item 学院官方群：初审通过后部分学院要求实名加群，四川大学生命科学学院要求初审通过者加入指定 QQ 群（认证格式「姓名-本科院校」），面试时间与分组名单在群内公布 \nref{S0245}。
  \item 官方联系方式：中山大学公开各招生单位官网入口与研究生招生单位通讯录 \nref{S0166}；东北师范大学提供《研究生招生学院（部）及咨询方式》，并说明专业、课程、导师信息等问题应咨询各招生院部 \nref{S0607}。
  \item 本校资源：向学院教务办或辅导员询问本校往届升入该院系的校友。
\end{itemize}
\nref{S0306}
\nref{S0304}
\nref{S0245}
\nref{S0166}
\nref{S0607}
\end{qa}

\iin{天津大学}
\ith{人脉渠道}
\ist{信息与工具}
\begin{qa}{S0380}{通过学长学姐打听消息时，怎么问才不冒犯又能拿到真话？}
核心是降低对方的回答成本与风险，并把问题问成可核实的事实，而不是评价。
\begin{itemize}
  \item 先给背景再提问：一句本科学校＋专业＋排名区间＋目标，对方才有判断适配度的依据。
  \item 问事实，不问态度：不问「这个导师好不好」，改问「他组里硕士几年毕业、周会频率、有没有横向项目」。接收方公开的程序事实可以当锚点，例如河海大学写明材料审核专家组不少于3人、面试专家组不少于5人、专家原则上应有副高以上职称 \nref{S0178}，中国矿业大学写明面试全程录音录像并对每位考生的作答情况进行现场记录 \nref{S0626}，你可以直接问「面试当天几个老师、有没有抽题」。
  \item 用公示数据做投射：把录取名单公示里公开的字段（如湖南大学计算机学院2027届推免名单含专业、姓名、学业成绩、综合排名、专业总人数、GPA、备注）\nref{S0131} 拿去问「你当年这批人本科大概什么层次」，比问「是不是只要985」更容易得到具体答案。
\end{itemize}
\nref{S0178}
\nref{S0626}
\nref{S0131}
\end{qa}

\iin{河海大学}
\ith{人脉渠道}
\ist{信息与工具}
\begin{qa}{S0381}{本科学校的校友资源在保研信息获取中能起多大作用？}
校友资源能提高命中率，但改不了规则层，要分清两件事。
规则层与你无关：所有推免生均享有依据招生政策自主选择报考研究生招生单位和专业的权利，推荐高校不得将报考本校作为遴选推免生的限制条件，也不得以其他任何形式限制推免生自主报考（教学厅〔2020〕12号）\nref{S0318}；教学〔2006〕14号也明确所有推免生均可报考其他招生单位，所有推免生名额均可用于向其他招生单位推荐 \nref{S0047}。所以「本校能不能推出去」本身不是问题。
信息层校友很关键：双非学生最缺的是「这个学院收过什么背景的人」。可用公开数据是拟录取名单公示（如湖南大学计算机学院2027届公示含专业、姓名、学业成绩、综合排名、专业总人数、备注）\nref{S0131}，把多年名单累积比对即可看出实际来源分布。另一路径是直接看接收方写明的申请条件——上海交通大学机械与动力工程学院把「毕业高校为『双一流』建设高校及建设学科或国家重点学科的优秀2027年应届本科生」写进申请条件 \nref{S0262}，这是可以核实的硬门槛。
\nref{S0318}
\nref{S0047}
\nref{S0131}
\nref{S0262}
\end{qa}

\iin{shpg（高校官网）}
\ith{人脉渠道}
\ist{信息与工具}
\begin{qa}{S0382}{怎么通过导师的关系网打听某位导师的真实情况？}
正规做法是顺着公开的学术关系找可核实的人，而不是托人递话。
可用路径：
\begin{itemize}
  \item 目标导师的合作者与同门：从课题组主页、论文合作者、实验室成员列表出发，找在读博士生或近期毕业生。学院层面的公开信息可缩小范围——北京大学历史学系设有研究生招生指南类栏目 \nref{S0200}，浙江大学信息与电子工程学院在夏令营通知中按研究所（方向）介绍主要研究方向与拟招营员专业 \nref{S0036}。
  \item 本校老师引荐：本科导师、任课老师与目标导师可能有学术交集，请他们以学术问题的方式引荐，比「帮忙问问」有效得多。
  \item 官方接触窗口：夏令营安排的师生面对面与实验室交流环节（天津大学软件学院明确「为营员搭建与导师、学长学姐面对面互动的平台」）\nref{S0306}。
\end{itemize}
要回避的：通过中介或付费机构打听导师「内部评价」。这类信息既无法核实，也可能在正式流程中给你留下不良印象。
\nref{S0200}
\nref{S0036}
\nref{S0306}
\end{qa}

\iin{北京大学}
\ith{人脉渠道}
\ist{信息与工具}
\begin{qa}{S0383}{找导师组里的在读博士生聊天，应该问哪些问题？}
按「能不能毕业、具体怎么过、值不值」三层问，全部问可核实的细节。
\begin{itemize}
  \item 时间与产出：组会频率与形式、从选题到投稿的典型周期、该组近三年硕博平均毕业年限。公开信息可做交叉验证：中国农业大学博士招生章程写明直博生学制5年、硕博连读博士阶段3年、最长在校学习年限为学制加2年 \nref{S0429}；河北大学写明直博生学制5年、最长在校学习年限不得超过8年 \nref{S0604}。
  \item 指导方式：导师多久见一次学生、具体课题由谁带（导师本人还是组内高年级博士）、是否有共同指导。
  \item 资源与风险：经费以纵向还是横向为主、有无企业联合培养、是否要求出国交流、论文署名惯例。
  \item 出路：近三年毕业生去向。可对照学校公开的学位与答辩信息，如上海财经大学研招网设「学位授予名单」「答辩公告」栏目 \nref{S0182}，东南大学材料科学与工程学院也会公示接收推免研究生的考核结果 \nref{S0585}。
\end{itemize}
\nref{S0429}
\nref{S0604}
\nref{S0182}
\nref{S0585}
\end{qa}

\iin{中国农业大学}
\ith{人脉渠道}
\ist{信息与工具}
\begin{qa}{S0401}{想找往年同背景的学生案例，有哪些渠道？}
优先找官方公开的结果数据，再找人证。
\begin{itemize}
  \item 拟录取名单公示：这是唯一系统性的官方来源，但过公示期可能下线，要在公示期内下载（南京理工大学2024年推免拟录取名单以附件形式公示）\nref{S0181}；学院层面也常单独公示，湖南大学计算机学院2027届推免名单公示含专业、姓名、学业成绩、综合排名、专业总人数、GPA、备注 \nref{S0131}。
  \item 学校「信息公开／历年数据」栏目：南京大学研招网设「信息公开—公示」「往年报考录取统计」\nref{S0505}；四川大学研招网设「推免生接受人数汇总表」「信息公开」「历年数据」\nref{S0472}；华东师范大学研招网设「历年硕士报录数据」\nref{S0611}。
  \item 夏令营与预推免的营员／复试名单公示，能看到同背景学生的分布。
  \item 人证渠道：夏令营的实验室交流与师生面对面环节（天津大学软件学院明确安排）\nref{S0306}；学院官方 QQ 群（四川大学生命科学学院要求初审通过者实名加群）\nref{S0245}。
\end{itemize}
\nref{S0181}
\nref{S0131}
\nref{S0505}
\nref{S0472}
\nref{S0611}
\nref{S0306}
\nref{S0245}
\end{qa}

\iin{南京理工大学}
\ith{人脉渠道}
\ist{信息与工具}
\begin{qa}{S0412}{邮件联系导师后没有回复，多久跟进一次比较合适？}
先明确一点：联系导师在多数学校不是录取的必要环节，回不回邮件不改变录取规则。决定结果的是系统流程——考生须在推免服务系统填报志愿，招生单位通过系统反馈是否同意复试，未通过系统确认接受复试通知的考生不能被招生单位录取 \nref{S0444}；上海交通大学写明以推免服务系统上发出的拟录取通知为准，考生在规定时间内予以答复确认，逾期则视为主动放弃 \nref{S0515}。
所以跟进策略不应以邮件为中心：
\begin{itemize}
  \item 时点：优先在学院组织复试前联系。学院考核时间由学院自定，华东师范大学心理与认知科学学院把综合考核安排在9月18—20日 \nref{S0295}，错过窗口再发邮件意义有限。
  \item 频率：未见统一规定。经验上同一导师跟进一次即可，不要在待录取阶段反复催促。
  \item 该盯的是官方通道：学校会通过推免服务系统、电话或电子邮件发送复试相关信息 \nref{S0067}。
\end{itemize}
\nref{S0444}
\nref{S0515}
\nref{S0295}
\nref{S0067}
\end{qa}

\iin{教育部学生服务与素质发展中心（学信网/研招网）}
\ith{联系导师}
\ist{信息与工具}
\begin{qa}{S0413}{给导师发邮件时，简历和成绩单应该怎么组织？}
按「学校要什么你给什么」的顺序组织，不要自己发明格式。
\begin{itemize}
  \item 先读目标校的材料清单，再决定附件。四川大学生命科学学院列出10项材料（含成绩单、排名证明、外语成绩、诚信承诺书，直博另需入学申请表）\nref{S0245}；上海交通大学要求申请材料包括个人基本信息、本科成绩单、外语水平证明、科研创新证明、重要奖项证书以及本人签名的诚信承诺书 \nref{S0515}；南京大学在接收推免办法下挂出「申请材料清单」「考生诚信复试承诺书」「专家推荐信」等附件模板 \nref{S0505}。
  \item 格式要求必须照做：多份材料常被要求合并为单个 PDF 并按规定命名。东北师范大学要求把多项材料合并为一个 PDF 上传至「佐证材料」栏，命名规则为「姓名-手机号码-材料名」\nref{S0607}；中国矿业大学要求电子申请材料按编号顺序装订成册并在学院规定截止日期前提交 \nref{S0626}。
  \item 成绩单要盖章：多数学校要求加盖推荐学校教务处公章的本科阶段学习成绩单 \nref{S0167}。
\end{itemize}
\nref{S0245}
\nref{S0515}
\nref{S0505}
\nref{S0607}
\nref{S0626}
\nref{S0167}
\end{qa}

\iin{四川大学}
\ith{联系导师}
\ist{信息与工具}
\section{信息渠道}
\begin{qa}{S0362}{招生简章和推免办法应该去哪些官方渠道查？}
优先按「接收方—推荐方—国家平台」三条线找，不要只在搜索引擎里碰运气。
\begin{itemize}
  \item 接收方：目标校研究生院/研究生招生信息网的「硕士招生—招生简章/招生专业目录」「推免」专栏；以及学院官网（接收实施细则、复试办法在这里）。武汉大学2027年公告写明「我校2027年接收推免生专业目录通过全国推免服务系统向社会公开，各培养单位接收推免生工作实施细则在本单位网站公布」\nref{S0347}。
  \item 推荐方：本校教务处/本科生院，推免办法、名额分配、加分细则在此。
  \item 国家平台：中国研究生招生信息网 yz.chsi.com.cn 的推免服务系统，招生单位专业目录与接收办法在其中公开。
\end{itemize}
依据：教学厅〔2013〕8号要求推荐单位和接收单位都要在本单位网站开辟专栏，向社会公布本单位推免生名额、推免办法，并对拟推荐名单、拟接收名单统一公示，公示时间不少于10个工作日 \nref{S0047}。这是学校必须做的事，你有权要求看到。
\nref{S0347}
\nref{S0047}
\end{qa}

\iin{武汉大学}
\ith{信息渠道}
\ist{信息与工具}
\begin{qa}{S0363}{学院官网、研究生院官网、推免服务系统信息不一致时以哪个为准？}
分层定优先级，不要笼统说「以官网为准」。
\begin{enumerate}
  \item 涉及你的资格与录取结果：以推免服务系统备案信息为准。研招网明确「最终推免生名单、录取名单和推免生学籍注册均以『推免服务系统』备案信息为准」\nref{S0464}；武汉大学同样写明「具有推免资格的考生名单以全国推免服务系统备案信息为准，未经全国推免服务系统备案的考生不得录取」\nref{S0347}。
  \item 时间与流程口径：系统开放等国家统一节点按教育部招考公告，学校自己的时间安排以学校通知为准。研招网对「学校页面显示的推免截止时间与网站上显示的不一致，以哪个为准」的回答是：本网站提示时间按教育部招考公告发布，学校的具体时间安排请与其联系咨询 \nref{S0523}。哈工大2027年提醒还要求系统内填写的报名志愿须与学院通知的拟选拔志愿一致，包括层次、专业、研究方向（直博生导师）\nref{S0067}。
  \item 学院细则与研究生院办法冲突：学院细则须报研究生院备案，冲突时以研究生院口径为准（中国矿业大学即要求细则报研究生院备案并在复试前公布）\nref{S0626}，遇到就直接打电话向研招办确认。
\end{enumerate}
\nref{S0464}
\nref{S0347}
\nref{S0523}
\nref{S0067}
\nref{S0626}
\end{qa}

\iin{教育部学生服务与素质发展中心（学信网/研招网）}
\ith{信息渠道}
\ist{信息与工具}
\begin{qa}{S0364}{往年推免拟录取名单在哪里查，公示期过了还能查到吗？}
三个可查位置，优先用官方公示。
\begin{itemize}
  \item 研究生院/研招网的「信息公开」「公示」「公告通知」栏：南京大学研招网设「信息公开—公示」与「往年报考录取统计」\nref{S0505}；江南大学研招网设「信息公示」\nref{S0629}；四川大学研招网设「信息公开」「历年数据」以及「推免生接受人数汇总表」\nref{S0472}；华东师范大学研招网设「历年硕士报录数据」\nref{S0611}。
  \item 学院官网通知公告：学院层面的名单常单独发，如南京理工大学研究生院发布《接收2024年推荐免试攻读研究生拟录取名单公示》，名单在附件里 \nref{S0181}。
  \item 推免服务系统：中国矿业大学写明「学校将通过网站主动公开接收推免生章程、推免生招生目录、复试办法、拟录取名单等全部信息，同时按教育部要求在『推免服务系统』公示公开」\nref{S0626}。
\end{itemize}
公示期过了还能不能查到：国家只规定公示时间不少于10个工作日 \nref{S0047}，并未规定公示页面须永久保留，因此过公示期链接可能下线或名单被撤。要往年数据，就得在公示期内自行下载附件与页面留档。 还要注意公示天数：国家统一要求不少于10个工作日 \nref{S0047}，但各校实际执行有差异（中国矿业大学为不少于7日 \nref{S0626}，湖南大学计算机学院2027届名单公示期为5天 \nref{S0131}），想留档就要在发布后立刻下载，不要按10个工作日给自己留缓冲。
\nref{S0505}
\nref{S0629}
\nref{S0472}
\nref{S0611}
\nref{S0181}
\nref{S0626}
\nref{S0047}
\nref{S0131}
\end{qa}

\iin{南京大学}
\ith{信息渠道}
\ist{信息与工具}
\begin{qa}{S0366}{各校研究生招生信息网的「信息公开」栏目一般包含什么？}
这是国家要求的固定动作，不是学校自发办的栏目。依据《2026年全国硕士研究生招生工作管理规定》第七十三条、七十五条：招生单位是信息公开责任主体，应按「谁公开、谁把关」「谁公开、谁解释」的原则在本单位招生网站公开年度招生信息；报名前公布本单位招生章程，其中应明确各学科专业的报考条件、拟安排招生计划、初试考试科目、复试方式及有关考核要求、录取程序和标准、学习方式、学制、学费标准、奖助办法等；复试前公布复试录取办法、各院系实施细则与进入复试的考生名单 \nref{S0051}。教学厅〔2019〕2号进一步要求提前公布招生章程、招生政策规定、招生专业目录、分专业招生计划、复试录取办法，所有拟录取名单由研究生招生管理部门统一公示不少于10个工作日，并在「全国硕士研究生招生信息公开平台」统一公开，未经公示的考生一律不得录取、不予学籍注册 \nref{S0431}。
实际栏目名可对照找：招生章程、招生专业目录、分专业招生计划、复试录取办法与实施细则、拟录取名单公示、往年报考录取统计、咨询与申诉渠道（武汉大学要求公开专业、招生计划、接收办法、拟录取名单及咨询申诉渠道）\nref{S0347}。
\nref{S0051}
\nref{S0431}
\nref{S0347}
\end{qa}

\iin{中华人民共和国教育部}
\ith{信息渠道}
\ist{信息与工具}
\begin{qa}{S0368}{推免服务系统里招生单位的推免章程在哪里找？}
入口在系统内，规则是「先看办法、再填志愿」。路径：登录 yz.chsi.com.cn/tm 进入推免服务系统，填报志愿前应认真阅读各招生单位发布的推免招生工作办法和专业目录。研招网《2025「推免服务系统」操作流程》把这一步写在「填报志愿」环节的前提里 \nref{S0444}；东华大学给考生的流程说明同样要求「填报志愿前应认真阅读推免招生工作办法和专业目录」\nref{S0419}。
系统内还能看到两类信息：一是招生单位在系统中发布的推免招生工作办法与专业目录；二是单位自设的补充材料栏，部分学校会在其中要求填写特定信息，例如南开大学物理科学学院要求推免生按系统附加材料的设置要求准确填写外语水平等信息、上传四六级成绩证书附件 \nref{S0291}。
注意系统与官网互补：天津大学写明「推免招生学科专业及拟招生人数，详见推免服务系统和研究生招生网公布的招生专业目录」\nref{S0520}，也就是两边都要看，不能只看系统里的简版。
\nref{S0444}
\nref{S0419}
\nref{S0291}
\nref{S0520}
\end{qa}

\iin{教育部学生服务与素质发展中心（学信网/研招网）}
\ith{信息渠道}
\ist{信息与工具}
\begin{qa}{S0370}{哪些渠道能拿到目标院校往年的推免真题、机试题？}
先做预期管理：官方渠道基本拿不到。福州大学在招生简章里明文写「我校一律不代购参考书，也不提供我校往年硕士研究生招生考试试卷」，并说明招生专业目录中所提供的书目为主要参考书目、课程（考试）大纲和参考书目录仅供参考，不作为出题范围 \nref{S0603}。所以「买真题」本身就不在国家制度保障范围内。
可用的合法替代路径：
\begin{itemize}
  \item 学校公开的复试录取实施细则里常有参考书目，如闽南师范大学文学院2026年复试细则按专业列出参考书目 \nref{S0638}；
  \item 学院官网的招生宣传与栏目，如上海交通大学机械与动力工程学院研究生教学网设「考研大纲」「如何选择导师」「常见问题」「往年统考复试分数线」等栏目 \nref{S0262}；
  \item 官方微信服务号与宣讲，如东北师范大学通过「东师研招」微信服务号推送招考信息 \nref{S0607}。
\end{itemize}
结论：真题只当线索，不当范围。
\nref{S0603}
\nref{S0638}
\nref{S0262}
\nref{S0607}
\end{qa}

\iin{福州大学}
\ith{信息渠道}
\ist{信息与工具}
\begin{qa}{S0397}{教育部的推免政策文件一般什么时候发布，在哪里查？}
两条主线，发布节奏不同。
\begin{enumerate}
  \item 年度招生工作管理规定：教育部每年9月印发下一年度《全国硕士研究生招生工作管理规定》，在教育部政府门户网站（moe.gov.cn）和中国研究生招生信息网（yz.chsi.com.cn）同步公布。2027年为教学〔2026〕2号，印发日期2026年9月21日 \nref{S0442}；2026年为教学〔2025〕2号，2025年9月15日印发 \nref{S0051}。
  \item 推免专项文件：教学〔2026〕2号第八十六条明确「推荐免试工作相关管理办法由教育部另文规定」\nref{S0442}。2027届推免中，多所高校文件引用的上位依据是《全国高等学校推荐优秀应届本科毕业生免试攻读研究生工作管理规定》（教学厅〔2026〕1号），例如云南师范大学 \nref{S0568}、哈尔滨师范大学 \nref{S0058}、中山大学 \nref{S0167} 均直接引用该文号。
\end{enumerate}
查法：在 moe.gov.cn 用「推免」「推荐免试」检索，在研招网「国家政策/教育部政策」栏目检索。
\nref{S0442}
\nref{S0051}
\nref{S0568}
\nref{S0058}
\nref{S0167}
\end{qa}

\iin{教育部学生服务与素质发展中心（学信网/研招网）}
\ith{信息渠道}
\ist{信息与工具}
\begin{qa}{S0406}{本校教务处的推免细则一般在哪里发布，怎么拿到往年版本？}
发布位置：本校教务处/本科生院官网的「通知公告」或「规章制度」栏，以及学校「信息公开」专栏。实例：湖南大学教务处发布《关于做好2027年湖南大学推荐优秀应届本科毕业生免试攻读硕士学位研究生工作的通知》\nref{S0540}；云南师范大学《关于做好2027年推荐优秀应届本科毕业生免试攻读研究生工作的通知》（云师发〔2026〕88号）发布在教务信息网 \nref{S0568}；内蒙古大学的实施办法挂在信息公开网 \nref{S0544}。
制度依据：教学厅〔2020〕12号要求推荐单位在本单位网站开辟专栏，公布本单位推免生名额、推免办法，并对拟推荐名单统一公示，公示时间不少于10个工作日 \nref{S0047}。公开是学校的义务。
往年版本怎么拿：一是校内信息公开专栏或制度汇编；二是教务处通知公告栏向前翻页；三是直接向教务处或学院教务办索取；四是请学长学姐提供他们当年的留档。
\nref{S0540}
\nref{S0568}
\nref{S0544}
\nref{S0047}
\end{qa}

\iin{湖南大学}
\ith{信息渠道}
\ist{信息与工具}
\section{信息甄别}
\begin{qa}{S0372}{怎么核实「某校只要985」这类说法的真伪？}
用「书面条件」和「实际名单」两条证据链去核，两步都能自己做完。
第一步，查书面条件。接收方若真有层次要求，会写进申请条件。上海交通大学机械与动力工程学院2027年预推免办法把「毕业高校为『双一流』建设高校及建设学科或国家重点学科的优秀2027年应届本科生」列为申请条件之一，并给出英语合格线 \nref{S0262}——这是可以引用的硬门槛。反之，若文件只写「具有推荐免试资格高校的应届本科毕业生」，那「只要985」的说法就没有文件依据，南昌大学基础医学院、中国地质大学（武汉）数理学院均只写这一条 \nref{S0001}\nref{S0324}。
第二步，查实际名单。把目标学院近三年拟录取名单的本科来源做频次统计（公示字段可参考湖南大学计算机学院2027届推免名单，含专业、姓名、学业成绩、综合排名、专业总人数、备注）\nref{S0131}。
另外，制度上不允许无依据歧视：招生单位不得按单位、行业、地域、学校层次类别等限定生源范围，也不得设置其他歧视性条件（教学厅〔2019〕2号）\nref{S0431}。
\nref{S0262}
\nref{S0001}
\nref{S0324}
\nref{S0131}
\nref{S0431}
\end{qa}

\iin{上海交通大学}
\ith{信息甄别}
\ist{信息与工具}
\begin{qa}{S0373}{保研群里的消息和官网信息，可信度怎么排序？}
排序原则是「谁对结果负责，谁的话可信」。
第一级，推免服务系统备案信息：最终推免生名单、录取名单和推免生学籍注册均以系统备案信息为准 \nref{S0464}。
第二级，招生单位官网文件：推荐单位和接收单位都要在本单位网站开辟专栏，公布推免生名额与推免办法，并统一公示拟推荐名单、拟接收名单，公示时间不少于10个工作日 \nref{S0047}；接收方的招生章程、专业目录、分专业招生计划、复试录取办法、拟录取名单属法定公开内容 \nref{S0431}。
第三级，学院官网细则：它是执行口径，但须报研究生院备案，冲突时以研究生院为准（中国矿业大学要求细则报研究生院备案并在复试前公布）\nref{S0626}。
第四级，微信服务号、短信、邮件等通知渠道（东北师范大学通过「东师研招」微信服务号推送招考信息）\nref{S0607}，以及系统内的复试与待录取通知 \nref{S0444}：这类是通知你办事，不构成规则解释。
第五级，群消息、经验帖、机构公众号。教育部明确对社会培训机构违规开展辅导培训或发布虚假招生宣传（广告）骗取钱财的，有关部门可责令限期改正、警告，有违法所得的没收，情节严重者责令停止招生、吊销办学许可证和营业执照 \nref{S0442}。
\nref{S0464}
\nref{S0047}
\nref{S0431}
\nref{S0626}
\nref{S0607}
\nref{S0444}
\nref{S0442}
\end{qa}

\iin{教育部学生服务与素质发展中心（学信网/研招网）}
\ith{信息甄别}
\ist{信息与工具}
\begin{qa}{S0374}{保研公众号的信息时效性和可靠性怎么判断？}
三个检验动作。
\begin{itemize}
  \item 看出处与日期：正规信息一定给原文链接、发布单位和发布日期。学校文件都有明确发布日，华东师范大学2027年预报名通知为2026-07-24、综合考核细则汇总为2026-09-15 \nref{S0611}\nref{S0613}，对不上就以官网为准。
  \item 看它是否改动关键数字：把公众号里的名额、截止时间、条件与原文对一遍。拟招生人数在招生专业目录里按专业逐个列出 \nref{S0414}；截止时间是硬信息，2027年10月18日截止填报志愿、10月20日关闭系统 \nref{S0464}。
  \item 看它是否借信息卖服务：教育部明确，对社会培训机构违规开展辅导培训或发布虚假招生宣传（广告）骗取钱财的，由当地教育行政部门或其他有关部门依据《民办教育促进法》《广告法》等责令限期改正并予警告，有违法所得的没收非法所得，情节严重的责令停止招生、吊销办学许可证和营业执照；同时严禁社会培训机构进入校园以张贴简章、广告等方式进行考研辅导培训宣传和组织活动 \nref{S0442}。
\end{itemize}
另外，正规高校会主动声明与中介无关，如东北师范大学写明从未委托任何中介机构或个人开展招生信息发布、宣传、报名等工作 \nref{S0607}。
\nref{S0611}
\nref{S0613}
\nref{S0414}
\nref{S0464}
\nref{S0442}
\nref{S0607}
\end{qa}

\iin{华东师范大学}
\ith{信息甄别}
\ist{信息与工具}
\begin{qa}{S0375}{知乎、小红书上的保研经验帖该怎么甄别？}
把经验帖当「假设生成器」，不当依据来源。
可用之处：帮你想到该查什么栏目。比如看到有人提到某院在系统外还有独立报名入口，就去核该院官网是否真有——中国矿业大学有研究生推免预报名系统 \nref{S0626}，南京航空航天大学有研究生报考服务系统 \nref{S0648}。
必须回官网核实的四类内容：一是政策与截止时间（推免服务系统2027年9月21日9:00开始填报志愿、10月18日截止填报、10月20日关闭）\nref{S0464}；二是名额与门槛（招生专业目录里的分专业拟招生人数 \nref{S0414}；申请条件里的层次与英语要求，如上海交通大学机械与动力工程学院 \nref{S0262}）；三是流程细节（中国矿业大学要求面试全程录音录像、材料审核专家组不少于3人 \nref{S0626}）；四是涉及付费服务的推荐（教学〔2026〕2号第八十一条对社会培训机构违规宣传有明确处置措施）\nref{S0442}。
判断技巧：帖子里缺年份、学院、专业、原文链接中任何一项，可信度都要大幅下调。
\nref{S0626}
\nref{S0648}
\nref{S0464}
\nref{S0414}
\nref{S0262}
\nref{S0442}
\end{qa}

\iin{中国矿业大学}
\ith{信息甄别}
\ist{信息与工具}
\begin{qa}{S0376}{怎么识别保研机构的营销性信息与焦虑贩卖？}
识别标准只有一条：它有没有在承诺制度上不可能的事。
制度底线一，不存在能操控录取结果的合法服务。《普通高等学校招生违规行为处理暂行办法》把高校「违规委托中介机构进行招生录取，或者以承诺录取为名向考生收取费用的」以及「发布违反国家规定的招生简章，或者进行虚假宣传、骗取钱财的」列为违规情形，由主管教育行政部门责令限期改正、给予警告或通报批评；情节严重的给予减少招生计划、暂停特殊类型招生试点项目或停止招生的处理；对直接负责的主管人员和其他直接责任人员视情节给予处分，涉嫌犯罪的移送司法机关 \nref{S0041}。
制度底线二，正规高校会主动与中介切割。东北师范大学写明「我校不举办或参与举办任何形式的考研辅导活动，从未委托任何中介机构或个人开展招生信息发布、宣传、报名等相关工作」\nref{S0607}；华东师范大学研招网设有「我校不举办任何考研辅导机构的公告」栏目 \nref{S0611}。
制度底线三，录取结果只走系统：推免生资格审核确认、报考、录取以及备案公开均通过推免服务系统进行，未经系统备案不得录取 \nref{S0347}。
因此凡出现「内部名额」「保过通道」「包录取」，均不可信；以「不报名就来不及」施压的焦虑话术同样不构成任何依据。
\nref{S0041}
\nref{S0607}
\nref{S0611}
\nref{S0347}
\end{qa}

\iin{中华人民共和国教育部}
\ith{信息甄别}
\ist{信息与工具}
\begin{qa}{S0377}{付费保研服务里哪些值得买、哪些是智商税？}
按「它是否提供了你无法自行获得的合法资源」来分，而不是按价格分。
不构成捷径的部分（无论收费多少都不改变结果）：录取必须通过推免服务系统完成——未通过系统确认接受复试通知的考生不能被招生单位录取，未通过系统确认接受待录取通知则录取不生效 \nref{S0444}；河北大学明确「未通过『推免服务系统』进行的拟录取无效」\nref{S0604}。因此任何「帮你操作系统」「内部推荐」的承诺都与制度冲突：《普通高等学校招生违规行为处理暂行办法》把高校违规委托中介机构进行招生录取、或以承诺录取为名向考生收取费用列为违规行为，可责令限期改正、警告、通报批评，情节严重的减少招生计划或暂停招生 \nref{S0041}。
可能是合法服务的部分（取决于你的时间成本）：材料格式整理与翻译润色、面试模拟、专业课知识梳理——它们作用于「给你自己做功」，不介入录取决策。判断标准是：这项服务是否把最终决定权完整地留给学校与系统。
另外，学校官方渠道免费且信息最全：招生院部咨询方式与官方微信服务号 \nref{S0607}、研招办与监察处的电话邮箱 \nref{S0626}。
\nref{S0444}
\nref{S0604}
\nref{S0041}
\nref{S0607}
\nref{S0626}
\end{qa}

\iin{教育部学生服务与素质发展中心（学信网/研招网）}
\ith{机构服务}
\ist{信息与工具}
\begin{qa}{S0378}{保研咨询机构的「保过协议」里有哪些坑？}
最大的坑是它承诺的事制度上做不到，退费条款写得再好也补不回你错过的窗口。
制度事实一，录取不走机构：推免生资格审核确认、报考、录取以及备案公开均通过推免服务系统进行，未经系统备案不得录取 \nref{S0347}；未通过系统确认接受待录取通知，则录取不生效 \nref{S0444}。
制度事实二，机构介入招生属违规：高校违规委托中介机构进行招生录取，或以承诺录取为名向考生收取费用，属《普通高等学校招生违规行为处理暂行办法》列明的违规情形，可责令限期改正、警告、通报批评，情节严重的减少招生计划、暂停特殊类型招生试点项目或停止招生 \nref{S0041}；高校也普遍声明未委托任何中介机构或个人开展招生信息发布、宣传、报名工作 \nref{S0607}。
制度事实三，弄虚作假的后果由你承担：对在推免过程中弄虚作假、有论文抄袭、虚报获奖或科研成果等学术不端行为的学生，一经查实取消推免资格，已入学的取消学籍，并由推荐高校所在地的省级教育招生考试机构按规定记入《国家教育考试考生诚信档案》\nref{S0318}；对违反国家招生规定取得入学资格或者学籍的，学校应当取消其学籍，以作弊、剽窃、抄袭等不正当手段获得的学历学位证书应予撤销 \nref{S0040}。
\nref{S0347}
\nref{S0444}
\nref{S0041}
\nref{S0607}
\nref{S0318}
\nref{S0040}
\end{qa}

\iin{武汉大学}
\ith{机构服务}
\ist{信息与工具}
\section{文件解读}
\begin{qa}{S0391}{官方文件里的「原则上」「一般应」这类表述怎么理解？}
它们不是模糊词，而是「有例外的规则」。读法是先找例外条件，再找谁有权批准例外。
在推免文件里常见三类用法：
\begin{itemize}
  \item 一般性比例／范围：教学〔2006〕14号写教育部批准设立研究生院的高等学校一般按应届本科毕业生数的15\%左右确定，未设立研究生院的「211工程」建设高校一般按5\%左右，其他高校一般按2\%（初次开展前3年一般按1\%），并写明教育部可根据形势对上述比例做适当调整 \nref{S0454}——「一般」意味着可调整，不构成你能主张的权利。
  \item 取值规则：教学厅〔2020〕12号写学生在某一方面中有多项加分情况时原则上只取一项，不得为仅符合单一或部分遴选指标的学生单列计划或破格推荐 \nref{S0318}——「原则上」后面紧跟的「不得」才是硬约束。
  \item 程序性要求：教学〔2026〕2号写原则上招生单位非全日制硕士研究生招收在职定向就业人员 \nref{S0051}。
\end{itemize}
判断方法：把「原则上／一般应」当默认值；要主张例外，就看文件有没有写明例外情形与批准机关。
\nref{S0454}
\nref{S0318}
\nref{S0051}
\end{qa}

\iin{教育部学生服务与素质发展中心（学信网/研招网）}
\ith{文件解读}
\ist{信息与工具}
\begin{qa}{S0392}{推免办法里「同等条件下优先」这类条款的实际含义是什么？}
它的含义是「不加分、只在并列时才起作用」，而且依附于文件中定义的特定情形。
典型用法一，学术专长：学生与直系亲属或学历、职称、职务明显高于本人者合作的科研成果、竞赛奖项等仅作为参考，不纳入学生本人推免遴选综合评价成绩计算体系，同等条件下可优先考虑 \nref{S0527}\nref{S0738}。典型用法二，申请条件：申请人在本科阶段公开发表过学术论文，或取得科研成果，或在国内外重大竞赛中获奖，可优先考虑 \nref{S0303}\nref{S0520}。典型用法三，专项计划：中共党员（含中共预备党员）同等条件下优先考虑 \nref{S0589}。
要点三条：一，它不改变排序分值，只在综合成绩完全相同时才成为比较依据；二，是否成立取决于文件对「条件」的定义（与谁合作、成果级别、身份），不符合定义连「优先」都不成立；三，「可优先考虑」比「同等条件下优先」语气更弱。
\nref{S0527}
\nref{S0738}
\nref{S0303}
\nref{S0520}
\nref{S0589}
\end{qa}

\iin{教育部学生服务与素质发展中心（学信网/研招网）}
\ith{文件解读}
\ist{信息与工具}
\begin{qa}{S0393}{怎么解读「接收推免生名额以实际下达为准」这句话？}
这句话通常对应两层含义，要分开看。
第一层，国家层面：专业目录里公布的是「拟」招生人数，可调整。中国农业大学博士招生章程写招生专业目录中公布的专业拟招人数仅供参考，最后录取以教育部实际下达招生计划为准 \nref{S0429}。天津大学写得更细：推免招生专业目录中所列各学科专业（研究方向）的拟招生人数均为根据以往生源情况所做的预估数，实际推免招生人数将根据录取情况，报经学校研究生招生委员会同意后进行调整 \nref{S0520}。中国矿业大学也写明各学院应根据目录公布的拟招收推免生人数开展复试录取，录取人数一般不超过拟招收推免生人数 \nref{S0626}。
第二层，名额回流：未招满的推免名额会用于统考——东华大学写明未招满的推免名额将用于统招 \nref{S0419}；武汉大学写明接收推免生工作结束后，培养单位将在本单位网站上公布可用于统考的硕士研究生分专业招生计划详情 \nref{S0347}。
国家规定的上限是：招生单位接收推免硕士生数量不得超过本单位全日制硕士研究生招生计划的50\% \nref{S0318}。
\nref{S0429}
\nref{S0520}
\nref{S0626}
\nref{S0419}
\nref{S0347}
\nref{S0318}
\end{qa}

\iin{中国农业大学}
\ith{文件解读}
\ist{信息与工具}
\begin{qa}{S0394}{怎么判断某校是「导师决定制」还是「学院委员会制」？}
不用听说法，去文件里找三处证据。
\begin{enumerate}
  \item 组织条款：看接收办法写的是谁「确定拟录取名单」。河海大学写明学校招生委员会审批拟录取名单，各单位成立研究生招生委员会（5—7人）并设秘书1名，另设材料审核专家组（不少于3人）与面试专家组（不少于5人）\nref{S0178}。出现「委员会＋专家组」结构，就是学院集体决策制。
  \item 程序条款：集体决策制通常伴随「集体研究、集体决策」「专家组现场独立评分」「面试全程录音录像」，如中国矿业大学要求面试全程录音录像、评分记录由研究生招生管理部门集中保管 \nref{S0626}，南京航空航天大学写明考核以综合面试为主、专家组一般不少于5人、面试全程录音录像 \nref{S0591}。
  \item 导师角色条款：导师话语权大的学校会正面写「需先联系导师确认接收意向」。中国地质大学（武汉）数理学院写申请考生需先跟博士生导师联系确认接收意向，并由博士生导师同意 \nref{S0324}；哈尔滨工业大学则要求系统内填写的报名志愿须与学院通知的拟选拔志愿一致，含研究方向（直博生导师）\nref{S0067}。
\end{enumerate}
\nref{S0178}
\nref{S0626}
\nref{S0591}
\nref{S0324}
\nref{S0067}
\end{qa}

\iin{河海大学}
\ith{文件解读}
\ist{信息与工具}
\begin{qa}{S0396}{保研时间节点以哪个官方文件为准？}
分两层：系统节点看教育部公告，校内节点看本校文件。
国家层：中国研究生招生信息网发布的推免服务系统时间安排是统一口径。2027年公开节点为：9月18日起注册填报个人资料并查看本人推免资格；9月21日9:00开始填报志愿；10月18日截止填报志愿；9月22日8:00开始确认复试通知；9月24日14:00开始确认待录取通知；10月20日关闭系统 \nref{S0464}。研招网操作流程把「填报志愿前应认真阅读各招生单位发布的推免招生工作办法和专业目录」写进第4步前提 \nref{S0444}。
学校层：各校自设节点更早也更细，河北大学公布的系统功能时间表为9月21日9:00开通报名、9月22日8:00开通复试通知、9月24日8:00发送待录取、10月18日关闭注册支付报名、10月20日关闭确认与发送功能 \nref{S0604}。
冲突时怎么办：研招网对「学校页面显示的推免截止时间与网站上显示的不一致」的回答是「本网站提示时间是按照教育部招考公告发布的，学校的具体时间安排请与其联系咨询」\nref{S0523}。
\nref{S0464}
\nref{S0444}
\nref{S0604}
\nref{S0523}
\end{qa}

\iin{教育部学生服务与素质发展中心（学信网/研招网）}
\ith{文件解读}
\ist{信息与工具}
\begin{qa}{S0398}{各校推免办法的修订历史能查到吗，怎么用？}
多半查得到，要用「制度文件」的方式查。
依据：推免办法属学校规章制度，修订时通常以「关于印发／修订通知」形式印发，正文写明审议通过日期与文号。南开大学《推荐优秀应届本科毕业生免试攻读研究生工作实施办法》以南党发〔2021〕63号印发，写明经2021年8月2日第二十四次党委常委会会议审议通过 \nref{S0561}；太原理工大学《推荐优秀应届本科毕业生免试攻读研究生工作管理办法（修订）》以校教〔2025〕6号发布 \nref{S0712}；广西大学《…工作实施办法（2023年修订）》以西大教〔2023〕83号印发 \nref{S0702}。这类版本通常在教务处或学校信息公开网的「规章制度」栏可查。
怎么用：把相邻两版并排，只看四处——学院推免细则应包含哪些内容、加分与学术专长认定范围、公示与复议条款、名额分配与排名课程确定程序。南开大学2021年版明确要求学院《推免细则》应包含学生报名、基本条件审核、综合素质评价、公示、复议等环节的具体操作方法和总体时间安排，并明确接受复议的时间、地点和联系方式 \nref{S0561}，这类条款一变，直接影响你的救济路径。
\nref{S0561}
\nref{S0712}
\nref{S0702}
\end{qa}

\iin{南开大学}
\ith{文件解读}
\ist{信息与工具}
\begin{qa}{S0407}{怎么确认加分政策的认定标准和材料清单？}
三步，缺一不可。
第一步，拿到本校今年的推免工作通知与学院推免细则。学院细则才是最终口径：南开大学要求学院《推免细则》明确推免生特殊学术专长审核办法，并保持相对稳定、确需调整须经学院推免生遴选工作领导小组批准 \nref{S0561}；西南交通大学要求各学院制定本学院推免工作实施细则，包含本学院加分细则，向学生公布并报学校推免生遴选工作领导小组备案 \nref{S0738}。
第二步，对照学术专长的硬边界。教学厅〔2020〕12号与各校落实口径基本一致：推免生学术专长原则上仅限本科阶段在核心期刊上以独立作者或第一作者发表的与学业相关论文，或作为主力成员参加与学业相关的国内权威科研竞赛（全国赛）并获得三等奖以上奖励；与直系亲属或学历、职称、职务明显高于本人者合作的成果不计入 \nref{S0527}\nref{S0738}；学生在某一方面中有多项加分时原则上只取一项 \nref{S0318}。
第三步，核对材料清单与形式。中南民族大学要求对竞赛奖项、论文、专利组织材料审核专家组（至少5名副教授以上专家）鉴定并组织公开答辩、全程录音录像，纸介质至少保留1年、电子扫描版至少保留5年、录音录像电子材料至少保留5年 \nref{S0594}；安徽师范大学还列出入伍建功、国际组织实习、志愿服务等专长对应分值 \nref{S0646}。
\nref{S0561}
\nref{S0738}
\nref{S0527}
\nref{S0318}
\nref{S0594}
\nref{S0646}
\end{qa}

\iin{南开大学}
\ith{文件解读}
\ist{信息与工具}
\section{信息渠道（续2）}
\begin{qa}{S0371}{考研真题对保研面试有参考价值吗？}
价值有限，而且容易把方向带偏。
\begin{itemize}
  \item 官方明确不提供：福州大学写明一律不代购参考书、不提供往年硕士研究生招生考试试卷，参考书目与考试大纲仅供参考，不作为出题范围 \nref{S0603}。也就是说，即便拿到真题，也不构成面试范围的任何承诺。
  \item 推免考核以综合面试为主：中国矿业大学写明综合考核包括材料审核、面试、思想政治素质和品德考核，还可根据学科专业需要增加笔试、实验考核、实践能力考核 \nref{S0626}；南京航空航天大学写明考核形式以综合面试为主，原则上按学院成立面试专家组，专家组一般不少于5人，面试全程录音录像 \nref{S0591}；华东师范大学心理与认知科学学院细则写明综合考核总成绩＝专业素质及综合能力面试×70\%＋英语能力×30\% \nref{S0295}。
\end{itemize}
可用之处只有一处：把真题当专业课基础的盘点清单，用来找自己哪门主干课有洞，而不是当押题。
\nref{S0603}
\nref{S0626}
\nref{S0591}
\nref{S0295}
\end{qa}

\iin{福州大学}
\ith{信息渠道}
\ist{信息与工具}
\begin{qa}{S0395}{怎么提前知道某校今年的政策变化（缩招、改直博）？}
靠纵向对比，不靠打听。三条可操作路径。
\begin{enumerate}
  \item 记住同校发布节奏，盯「陆续更新」的汇总页。多数学校在6—9月分批发布：华东师范大学2026年7月24日发布各院系预报名通知（标注「陆续更新」），9月15日发布综合考核实施细则汇总 \nref{S0611}\nref{S0613}；中国矿业大学2026年9月8日发布接收通知 \nref{S0626}；北京航空航天大学设「各学院接收推免攻读2027年研究生相关通知（2026年9月26日更新）」页 \nref{S0065}。
  \item 逐条对比同校上下两年文本，重点看四处：招生专业目录中的拟招生人数（数字变化即缩招/扩招）\nref{S0414}、招生类型是否新增直博、考核方式与成绩构成、材料清单增减。
  \item 读措辞变化。南开大学物理科学学院2027年通知写「后期是否再次开放报名，将视是否已完成招生计划待定」\nref{S0291}，天津大学写明目录中拟招生人数为按以往生源情况做的预估数、实际会根据录取情况调整 \nref{S0520}——这类表述本身就是信号。
\end{enumerate}
\nref{S0611}
\nref{S0613}
\nref{S0626}
\nref{S0065}
\nref{S0414}
\nref{S0291}
\nref{S0520}
\end{qa}

\iin{华东师范大学}
\ith{信息渠道}
\ist{信息与工具}
\begin{qa}{S0399}{怎么核实一所学校是否具有接收推免生的资格？}
先分清两个不同的名单：接收推免生的资格，与开展推荐工作的资格，不是一件事。
推荐资格的产生程序：教学〔2006〕14号规定，高等学校认为本校符合条件的，可向学校所在地省级研究生招生主管部门提出申请，省级主管部门核实并遴选后转报教育部，教育部组织推免生工作专家委员会评议后向社会公示，公示后无异议的，由教育部通知并公布开展推荐工作的高等学校名单 \nref{S0454}。
接收方与学科的关系：《2027年全国硕士研究生招生工作管理规定》写明，开展单独考试招生、推荐免试招生中推荐工作的单位应具备相应资格；凡按规定可接受应届本科毕业生报考的学科（类别）、专业（领域）均可接收推免生，但全日制学科专业不得只接收推免生 \nref{S0442}。
实操核实两步：一是在推免服务系统里检索该单位当年是否发布推免章程与专业目录；二是看该单位研究生院官网是否有当年的接收推免生办法与专业目录（武汉大学、中国矿业大学均在官网公开发布）\nref{S0347}\nref{S0626}。
\nref{S0454}
\nref{S0442}
\nref{S0347}
\nref{S0626}
\end{qa}

\iin{教育部学生服务与素质发展中心（学信网/研招网）}
\ith{信息渠道}
\ist{信息与工具}
\begin{qa}{S0400}{怎么查一所本科学校是否具有推免资格、名额有多少？}
两个问题分开查。
推免资格（学校有没有）：教育部按程序审核并向社会公示、公布开展推荐工作的高等学校名单 \nref{S0454}；新增名单另行备案公布，例如2025年教育部办公厅《关于对新增推免资格高校予以备案的通知》备案了北京电子科技学院等67所新增推免资格高校 \nref{S0640}。另有一个经验判据：每年发布推免工作通知的学校才是有名额的学校。
名额（有多少）：教学〔2006〕14号给出基准比例——教育部批准设立研究生院的高等学校一般按应届本科毕业生数的15\%左右确定；未设立研究生院的「211工程」建设高等学校一般按5\%左右；其他高等学校一般按2\%确定，其中初次开展推荐工作的前3年每年一般按1\%确定；人文、理科等人才培养基地的高校，按基地班招生人数的50\%左右单独增加名额 \nref{S0454}。该规定同时明确：高等学校不得将推免生名额跨学校使用；被确定的推免生如果无招生单位接收或本人放弃，推免生资格作废，名额不得转让 \nref{S0454}。
\nref{S0454}
\nref{S0640}
\end{qa}

\iin{教育部学生服务与素质发展中心（学信网/研招网）}
\ith{信息渠道}
\ist{信息与工具}
\begin{qa}{S0405}{怎么利用学校的就业指导中心、学业指导中心获取保研信息？}
先说结论：未见规定要求就业指导中心承担推免信息服务，它不在推免工作的信息责任链条上。推免信息公开的责任主体是推荐单位与接收单位——教学厅〔2020〕12号要求推荐单位和接收单位都在本单位网站开辟专栏，公布本单位推免生名额、推免办法，并统一公示拟推免名单、拟接收名单，公示时间不少于10个工作日 \nref{S0047}；接收单位的招生信息公开要求见教学〔2025〕2号第七十三至七十五条 \nref{S0051}。所以不要指望就业指导中心替你发布推免规则。
能实际用上的三处：一是校内部门链接里，就业指导中心常与研究生院、教务处并列（如南京理工大学研究生院页脚、江南大学研招网部门链接）\nref{S0180}\nref{S0629}，可借此找到校级咨询入口；二是就业指导中心办的升学与生涯讲座常请研究生院老师讲政策，现场提问比邮件有效；三是学院层面的推免说明会通常由教务办或研究生院组织。
\nref{S0047}
\nref{S0051}
\nref{S0180}
\nref{S0629}
\end{qa}

\iin{中华人民共和国教育部}
\ith{信息渠道}
\ist{信息与工具}
\begin{qa}{S0419}{有哪些非官方但信息质量较高的保研信息源值得长期跟踪？}
判断标准比渠道清单更重要：看它是否指向可核验的一手文件。
优先长期跟踪这四类：
\begin{itemize}
  \item 目标高校研究生院官方微信服务号，如东北师范大学「东师研招」微信服务号推送招考信息 \nref{S0607}；
  \item 学校研招网的「最新通知/公告通知」栏，如江南大学各学院接收推免细则汇总页、四川大学研招网的公告通知与信息公开 \nref{S0629}\nref{S0472}、华东师范大学的预报名与细则汇总（标注陆续更新）\nref{S0611}\nref{S0613}；
  \item 学院官网通知公告，接收细则、复试办法、拟录取公示的第一手发布地（武汉大学要求各培养单位在本单位网站公布实施细则并公开拟录取名单与咨询申诉渠道）\nref{S0347}；
  \item 中国研究生招生信息网的国家政策与推免服务栏目，用来看教育部原文 \nref{S0442}。
\end{itemize}
非官方渠道（群、公众号、论坛）的定位只有一个：提示线索。教育部明确，对社会培训机构违规开展辅导培训或发布虚假招生宣传（广告）骗取钱财的，由当地教育行政部门或其他有关部门依《民办教育促进法》《广告法》责令限期改正、予以警告，有违法所得的没收，情节严重的责令停止招生、吊销办学许可证和营业执照 \nref{S0442}。
\nref{S0607}
\nref{S0629}
\nref{S0472}
\nref{S0611}
\nref{S0613}
\nref{S0347}
\nref{S0442}
\end{qa}

\iin{东北师范大学}
\ith{信息渠道}
\ist{信息与工具}
\section{数据分析}
\begin{qa}{S0365}{怎么查目标院校的报录比和推免生占比？}
严格意义的「推免报录比」官方不发布，能查到的是三类替代数据，要自己拼。
\begin{itemize}
  \item 名额端：招生专业目录里的分专业拟招生人数。华北电力大学2027年推荐免试研究生招生专业目录按学院、专业代码、学位类型、专业名称逐项列出拟招生人数与研究方向 \nref{S0414}；天津大学明确目录中的拟招生人数为根据以往生源情况所做的预估数 \nref{S0520}。
  \item 实际录取端：拟录取名单公示里的实际人数，如南京理工大学2024年推免拟录取名单以附件公示 \nref{S0181}。把多年名单累积成时间序列，就能看出波动。
  \item 结构约束：招生单位接收推免硕士生数量不得超过本单位全日制硕士研究生招生计划的50\%，各全日制专业必须留出一定比例用于考试招生 \nref{S0318}；接收结束后学校会公布可用于统考的分专业招生计划（武汉大学）\nref{S0347}，用总计划减去它即可反推推免占用。
  \item 学校历史数据栏目：华东师范大学研招网设「历年硕士报录数据」\nref{S0611}，南京大学研招网设「往年报考录取统计」\nref{S0505}。
\end{itemize}
\nref{S0414}
\nref{S0520}
\nref{S0181}
\nref{S0318}
\nref{S0347}
\nref{S0611}
\nref{S0505}
\end{qa}

\iin{华北电力大学}
\ith{数据分析}
\ist{信息与工具}
\begin{qa}{S0367}{怎么从拟录取名单反推某校的招生偏好？}
名单里能读出的偏好只有三类，其余都是猜测。
\begin{itemize}
  \item 本科来源分布：这是唯一能直接读出的偏好，把多年公示名单的本科院校做频次统计。
  \item 排序规则：公示的字段与细则里的成绩构成能反推「什么更重要」。武汉大学要求各培养单位在本单位官方网站公开拟录取推免生名单（含姓名、复试各项成绩等）\nref{S0347}；华东师范大学心理与认知科学学院写明综合考核总成绩＝专业素质及综合能力面试×70\%＋英语能力×30\%，说明该项目面试权重远高于英语 \nref{S0295}；中国矿业大学写明材料审核专家组不少于3人、面试专家组不少于5人 \nref{S0626}。
  \item 名额结构：招生专业目录里的分方向拟招生人数决定哪个方向更容易进 \nref{S0414}，天津大学提示目录数字为预估数、会根据录取情况调整 \nref{S0520}。
\end{itemize}
读不出来的：导师个人偏好与是否存在非公开筛选。对无法从公示字段推出的结论，不要当成事实。
\nref{S0347}
\nref{S0295}
\nref{S0626}
\nref{S0414}
\nref{S0520}
\end{qa}

\iin{武汉大学}
\ith{数据分析}
\ist{信息与工具}
\begin{qa}{S0369}{怎么查到某专业今年到底招多少推免生？}
唯一权威口径是当年的招生专业目录，但要区分「拟」与「实际」。
\begin{itemize}
  \item 找法：目标校研究生招生信息网的「招生专业目录／推免招生专业目录」，以及推免服务系统内公开的招生单位专业目录。天津大学明确「推免招生学科专业及拟招生人数，详见推免服务系统和研究生招生网公布的招生专业目录」\nref{S0520}；中国药科大学写明「各学科招收推免生人数详见《中国药科大学2027年拟招收推荐免试研究生（含直博生）专业目录》」\nref{S0418}。
  \item 目录里看什么：按专业（方向）逐项列的拟招生人数，华北电力大学2027年目录精确到研究方向层级 \nref{S0414}。
  \item 口径提醒：目录数字是预估数，实际会按录取情况调整（天津大学）\nref{S0520}；中国农业大学写明专业目录中拟招人数仅供参考，最后录取以教育部实际下达招生计划为准 \nref{S0429}。
  \item 上限约束：接收推免硕士生数量不得超过本单位全日制硕士研究生招生计划的50\% \nref{S0318}。
\end{itemize}
\nref{S0520}
\nref{S0418}
\nref{S0414}
\nref{S0429}
\nref{S0318}
\end{qa}

\iin{天津大学}
\ith{数据分析}
\ist{信息与工具}
\begin{qa}{S0402}{怎么用录取名单数据评估自己的上岸概率？}
精确概率算不出来，但「够不够格进复试」可以估。三步。
第一步，摸分母：目标专业当年公布的拟招收推免生人数（见招生专业目录，如华北电力大学2027年目录按专业与研究方向列出拟招生人数）\nref{S0414}，以及接收推免硕士生不得超过本单位全日制硕士招生计划50\%的上限约束 \nref{S0318}。天津大学提示目录中拟招生人数为按以往生源情况做的预估数、实际会根据录取情况调整 \nref{S0520}。
第二步，摸分子分布：近三年拟录取名单的本科来源与成绩区间。可参考的公开字段如湖南大学计算机学院2027届推免名单公示含专业、姓名、学业成绩、所在专业学业成绩排名、综合成绩、综合排名、专业总人数、GPA、备注 \nref{S0131}；再看复试成绩构成，如华东师范大学心理与认知科学学院综合考核总成绩＝专业素质及综合能力面试×70\%＋英语能力×30\% \nref{S0295}。
第三步，给自己定位：把「我的本科层次＋排名区间」在近三年名单中出现的频次当作经验概率，而不是精确概率。若接收方申请条件写明层次要求，那是硬门槛，不是概率问题。
\nref{S0414}
\nref{S0318}
\nref{S0520}
\nref{S0131}
\nref{S0295}
\end{qa}

\iin{华北电力大学}
\ith{数据分析}
\ist{信息与工具}
\begin{qa}{S0418}{怎么用公开数据判断某校某专业的推免竞争烈度？}
用五个公开指标拼出烈度画像，不需要内部信息。
\begin{enumerate}
  \item 名额规模：招生专业目录里该专业（方向）的拟招生推免人数，华北电力大学2027年目录精确到研究方向 \nref{S0414}。个位数名额的方向，波动远大于趋势。
  \item 名额变化：与上年目录对比，缩招是最强信号。天津大学提示目录中的拟招生人数是按以往生源情况做的预估数、会根据录取情况调整 \nref{S0520}。
  \item 门槛表述：申请条件是否写死本科层次或英语线。上海交通大学机械与动力工程学院要求毕业高校为「双一流」建设高校及建设学科或国家重点学科，并要求 CET6≥425 或 TOEFL≥90 或 IELTS≥6.0 \nref{S0262}——一旦写明，就不是概率问题。
  \item 复试结构：河海大学要求材料审核专家组不少于3人、面试专家组不少于5人 \nref{S0178}；华东师范大学心理与认知科学学院综合考核中面试占70\% \nref{S0295}。面试权重高意味着名次靠前也可能被翻盘。
  \item 时间节奏：东华大学在流程说明后附「各学院推免生录取情况（实时更新）」表，逐院标注「各专业继续招收」或「推免名额已满」\nref{S0419}——这是最直接的实时烈度指标。
\end{enumerate}
\nref{S0414}
\nref{S0520}
\nref{S0262}
\nref{S0178}
\nref{S0295}
\nref{S0419}
\end{qa}

\iin{华北电力大学}
\ith{数据分析}
\ist{信息与工具}
\begin{qa}{S0420}{怎么判断一份往年名单是统考录取还是推免录取？}
看发布时间、标题与字段三条线索，不要靠人数猜。
\begin{itemize}
  \item 时间：推免录取在9—10月完成（2027年为9月21日填报志愿、10月18日截止填报志愿、10月20日关闭系统）\nref{S0464}；东华大学称2027年推免生拟录取名单预计10月中旬公示 \nref{S0419}。统考复试录取在次年3—4月（如福州大学写明复试预计安排在2026年3月下旬至4月中下旬）\nref{S0603}。发布月份是最快的判据。
  \item 标题与发布主体：推免名单标题多含「推免」「推荐免试」「接收」，且常由研究生院统一公示，如南京理工大学《接收2024年推荐免试攻读研究生拟录取名单公示》\nref{S0181}；统考名单标题含「复试录取」「拟录取」并挂在复试办法下。
  \item 字段：推免名单常含本科院校、推免成绩排名、复试各项成绩等（武汉大学要求公开拟录取推免生名单，含姓名、复试各项成绩等）\nref{S0347}；统考名单字段含考生编号、初试各科成绩（招生单位须公布进入复试考生名单，包括考生姓名、考生编号、初试各科成绩等）\nref{S0051}。
\end{itemize}
另外，推免录取必须经推免服务系统备案，未经系统备案不得录取 \nref{S0347}。
\nref{S0464}
\nref{S0419}
\nref{S0603}
\nref{S0181}
\nref{S0347}
\nref{S0051}
\end{qa}

\iin{教育部学生服务与素质发展中心（学信网/研招网）}
\ith{数据分析}
\ist{信息与工具}
\section{导师调研}
\begin{qa}{S0384}{通过领英、学术主页、谷歌学术调研导师，重点看什么？}
看四组可核验信号，并且优先用中文官方来源交叉验证。
\begin{itemize}
  \item 发文节奏与署名位置：近三年是否有稳定产出；合作者列表能反映组内规模与流动。
  \item 课题来源与方向漂移：近五年研究主题是收敛还是发散。
  \item 招生痕迹：学校官网的导师信息与院系方向介绍。浙江大学信息与电子工程学院在夏令营通知中按研究所（方向）介绍主要研究方向与拟招营员专业 \nref{S0036}；北京大学历史学系设有研究生招生指南类栏目 \nref{S0200}。
  \item 学生出口：结合学校公开的学位与答辩信息（上海财经大学研招网设「导师队伍」「学位授予名单」「答辩公告」等栏目）\nref{S0182}。
\end{itemize}
关键提醒：项目与经费的活跃度在国内不通过领英或谷歌学术体现，这些平台的信息也常滞后一两年。判断「现在是否还活跃」，最有效的是看学校官网最近一次发布的导师名单或招生方向目录。教育部办公厅亦要求进一步规范和加强研究生培养管理，导师履职属被管理事项 \nref{S0432}。
\nref{S0036}
\nref{S0200}
\nref{S0182}
\nref{S0432}
\end{qa}

\iin{浙江大学}
\ith{导师调研}
\ist{信息与工具}
\begin{qa}{S0385}{怎么判断一位导师近年是否还活跃（招生、发文、经费）？}
用「是否出现在今年的招生文件里」作第一判据，比数论文更直接。
\begin{itemize}
  \item 招生层面：看目标院系当年夏令营／预推免通知与专业目录中列出的研究方向与导师信息。浙江大学信息与电子工程学院按研究所列出主要研究方向与拟招营员专业 \nref{S0036}；天津大学化工学院夏令营通知把研究方向与活动内容一并公布 \nref{S0304}；中国地质大学（武汉）数理学院在接收办法中写明申请直博的考生需先跟博士生导师联系确认接收意向，说明导师是否参与招生会在办法里体现 \nref{S0324}。
  \item 培养层面：看院系公示的答辩公告与学位授予名单（上海财经大学研招网设「学位授予名单」「答辩公告」栏目）\nref{S0182}；中国农业大学博士招生章程会写明招生方式与规模安排 \nref{S0429}。
  \item 制度背景：教育部办公厅要求进一步规范和加强研究生培养管理，导师的招生资格与履职是被纳入管理的事项 \nref{S0432}。
\end{itemize}
论文数量只是间接指标，不同学科发文节奏差异极大，不能单独作为判断依据。
\nref{S0036}
\nref{S0304}
\nref{S0324}
\nref{S0182}
\nref{S0429}
\nref{S0432}
\end{qa}

\iin{浙江大学}
\ith{导师调研}
\ist{信息与工具}
\begin{qa}{S0386}{怎么判断一个课题组的毕业生去向是否被美化？}
用结构化数据替代故事。
\begin{itemize}
  \item 找原始名单而不是宣传稿：学校公开的学位授予与答辩信息属原始记录（上海财经大学研招网设「学位授予名单」「答辩公告」栏目）\nref{S0182}；院系公示的考核结果也属原始记录，如东南大学材料科学与工程学院公示2027年接收推荐免试研究生综合能力考核结果（第二批）\nref{S0585}。
  \item 算两个数：该课题组近三年毕业生人数（分母）与进入高校、科研院所、头部企业的比例（分子）。只看到「有人去了名校」这类孤例、看不到同届其他人去向的，一律按被美化处理。
  \item 交叉验证毕业节奏：中国农业大学博士招生章程写明直博生学制5年、硕博连读博士阶段学制3年、最长在校学习年限为学制加2年 \nref{S0429}；河北大学写明直博生学制5年、在校最长学习年限不得超过8年 \nref{S0604}。把公开学制与实际毕业年份对齐，能看出是否普遍延期。
  \item 找同期在读生核对：学院夏令营安排的师生面对面与实验室交流环节是正当的核对窗口 \nref{S0306}。
\end{itemize}
\nref{S0182}
\nref{S0585}
\nref{S0429}
\nref{S0604}
\nref{S0306}
\end{qa}

\iin{上海财经大学}
\ith{导师调研}
\ist{信息与工具}
\begin{qa}{S0411}{怎么查目标院校导师的招生名额是否已满？}
三个层次，从最直接到最间接。
\begin{itemize}
  \item 第一层，看学校实时公布的录取情况。东华大学在推免流程说明后附「各学院推免生录取情况（实时更新）」表，逐院标注「各专业继续招收，欢迎报考！」或「推免名额已满」，涵盖纺织学院、材料科学与工程学院、机械工程学院、数学与统计学院、人文学院等 \nref{S0419}。这是公开且最直接的判据。
  \item 第二层，看目录与办法里的名额口径：分专业（方向）的拟招生人数在招生专业目录中列出 \nref{S0414}；中国矿业大学写明各学院应根据目录公布的拟招收推免生人数开展复试录取，录取人数一般不超过拟招收推免生人数 \nref{S0626}；天津大学提示目录中的拟招生人数为预估数、实际会根据录取情况调整 \nref{S0520}。
  \item 第三层，直接问：院系研招办或导师本人。中国地质大学（武汉）数理学院要求直博申请者先跟博士生导师联系确认接收意向 \nref{S0324}；东北师范大学把导师信息类问题指向各招生院部咨询 \nref{S0607}。
\end{itemize}
\nref{S0419}
\nref{S0414}
\nref{S0626}
\nref{S0520}
\nref{S0324}
\nref{S0607}
\end{qa}

\iin{yjszs（高校官网）}
\ith{导师调研}
\ist{信息与工具}
\section{系统操作}
\begin{qa}{S0415}{推免系统操作有哪些常见错误需要提前知道？}
六个高频错误，每一个都有制度后果。
\begin{enumerate}
  \item 错过确认时限：复试通知与待录取通知都必须在招生单位规定时间内回复，未按时回复学院可取消（东华大学写明「若未按时回复，学院可取消复试通知／待录取通知」）\nref{S0419}；国防科技大学写明考生须在规定时间内在系统中同意待录取，否则视为放弃录取资格 \nref{S0650}。
  \item 未在系统内确认接受复试：未通过系统确认接受复试通知的考生不能被招生单位录取 \nref{S0444}。
  \item 接受待录取后想反悔：待录取通知可拒绝多个但只能接受一个，接受后不能再填报其他志愿，也不能接受其他复试与待录取通知 \nref{S0444}。
  \item 志愿填得与预选拔不一致：哈工大2027年提醒要求系统内填写的报名志愿须与学院通知的拟选拔志愿一致，包括层次（硕士／直博）、专业、研究方向（直博生导师）\nref{S0067}。
  \item 以为只在系统报名就够：多数学校另有预报名系统，东北师范大学写明已在学校系统报名的申请者仍须在教育部推免服务系统中报名，未经过系统接收的推免生一律无效 \nref{S0607}；河北大学写明未通过推免服务系统进行的拟录取无效 \nref{S0604}。
  \item 信息填报不实：因填写信息虚假、错误造成的后果由本人承担 \nref{S0444}。
\end{enumerate}
\nref{S0419}
\nref{S0650}
\nref{S0444}
\nref{S0067}
\nref{S0607}
\nref{S0604}
\end{qa}

\iin{yjszs（高校官网）}
\ith{系统操作}
\ist{信息与工具}
\begin{qa}{S0416}{推免系统的通知机制是怎样的，会不会漏掉？}
系统会发，但责任在你。机制是：推免生提交志愿后，招生单位通过推免服务系统反馈是否同意参加复试；收到复试通知后须在招生单位规定时间内通过系统回复，否则招生单位可取消复试通知；通过复试后招生单位通过系统发送待录取通知，考生须在规定时间内通过系统答复，否则招生单位可取消待录取通知；未通过系统确认接受待录取通知的，录取不生效 \nref{S0444}。
并行的通知渠道还有四类：电话或电子邮件（哈工大写明将通过全国推免服务系统、电话或电子邮件等方式向考生发送复试相关信息）\nref{S0067}；官方微信服务号推送（东北师范大学「东师研招」）\nref{S0607}；学院群或邮件（四川大学生命科学学院要求初审通过者实名加群，面试安排群内公布并请考生保持通讯畅通）\nref{S0245}；邮件为主的通知（华东师范大学心理与认知科学学院写明具体安排以邮件通知为准，请推免生注意查收邮件）\nref{S0295}。
所以会不会漏掉，取决于你有没有同时管好系统站内、邮箱、手机、公众号四个入口。
\nref{S0444}
\nref{S0067}
\nref{S0607}
\nref{S0245}
\nref{S0295}
\end{qa}

\iin{教育部学生服务与素质发展中心（学信网/研招网）}
\ith{系统操作}
\ist{信息与工具}

\begin{qa}{N0008}{推免系统里的政策、专业目录和院校公告，应该按什么顺序核对？}
先用教育部和省级教育招生考试机构公告确认全国时间与基本程序；再看目标招生单位当年的接收办法、专业目录和学院实施细则；最后以推免服务系统当前页面里的志愿、通知和备案状态完成操作。若同一事项出现冲突，记录文件标题、发布日期、适用年度和原文截图，联系发布单位确认并留存书面答复。系统中的公开材料由推荐高校和招生单位提供，平台展示本身不替代对原发布文件的核验。\nref{S0761}
\qscope{全国推免服务系统2027届说明；具体招生条件以对应学校当年正式文件为准。}
\end{qa}

\begin{qa}{N0009}{2027届推免政策页面更新后，怎样确认自己看到的不是旧缓存？}
从研招网推免服务系统和目标单位官网直接进入，不只依赖搜索引擎摘要、群聊转发或浏览器历史记录。记录访问日期、页面标题和发布/更新时间，下载正式附件并检查文件首页、落款日期和适用届别；若页面提示内容调整，重新下载附件并对照变化。关键志愿和待录取操作前再次刷新官方页面，避免按截图中的旧时间点操作。\nref{S0761}
\qscope{全国推免服务系统页面核对于2026年10月12日；网页和附件后续可能更新。}

核对时把网页内容与原始附件一起留存：记录网址、页面标题、发布日期/更新时间和查看日期，下载并保存带落款的PDF或公告附件；重要变化可记下新旧版差异和变更时间。刷新仍显示旧内容时从主管部门或招生单位官网首页重新进入，或换浏览器核验，并向发布单位确认。截图适合说明当时页面状态，不能替代正式文件；后续操作应按最新有效版本执行。
\end{qa}

\begin{qa}{N0010}{研招网公开的推免招生办法是谁提供的？出现错误应找哪个部门？}
研招网说明，推免系统中的推荐办法、推免生名单、招生办法和拟录取名单等由推荐高校或招生单位提供并按规定公示。资格推荐信息有误，先联系本科推荐高校教务部门或推免工作组；接收、复试和拟录取信息有误，联系目标招生单位研招办或学院。准备本人信息、页面截图和原始公示链接，按官方咨询申诉渠道反馈；涉及个人数据时避免在公开社交平台张贴证件或完整成绩单。\nref{S0761}
\qscope{全国推免服务系统2027届说明；咨询申诉路径按对应高校和省级机构公告。}
\end{qa}

\begin{qa}{N0011}{查到同专业其他推免生的录取公示后，可以联系对方打听未公开的申请材料吗？}
公示公开的是招生单位依法依规发布的录取信息，不代表其他申请人的成绩单、联系方式、材料或面试细节可随意索取和传播。用公开名单时只记录与决策相关的汇总信息，例如专业、层次和公示日期；如需了解培养体验，可自愿联系校友并尊重对方不回答的选择。不要转载身份证号、电话、个人邮箱或未公开材料。\nref{S0761}
\qscope{全国推免服务系统2027届说明仅描述录取公示范围；个人信息处理遵守适用法律及学校要求。}

联系前先把问题限定为自愿分享的经验，例如准备节奏、公开考核形式或入学培养体验；不要要求对方提供成绩单、证件、导师邮件、面试评分或未公开的个人材料，也不要把公示名单截图转发到公开群。若只是评估竞争情况，优先用招生单位公开的专业、层次、类别和年份做汇总，不从零散个人信息推断录取偏好。对方拒绝或不回复时应停止追问。
\end{qa}
\begin{qa}{S0417}{手机号、邮箱变更对推免流程有什么影响？}
会影响，因为系统和学校都用它推送关键通知，但制度上并不保证「换了就自动同步」。
\begin{itemize}
  \item 系统侧：注册学信网账号后须牢记注册手机号、邮箱、证件号码、密码等信息 \nref{S0523}；账号与学籍信息绑定，登录学信档案后可核实本人学籍学历信息 \nref{S0182}。手机号是账号找回与接收通知的主要通道，变更前先在系统内更新并完成验证。
  \item 学校侧：学校会通过推免服务系统、电话或电子邮件发送复试相关信息 \nref{S0067}；青岛大学提示请考生务必保持通讯畅通 \nref{S0693}；华东师范大学心理与认知科学学院写明具体安排以邮件通知为准、请推免生注意查收邮件 \nref{S0295}；东北师范大学提示因考生个人疏忽等原因造成的一切后果由考生本人承担 \nref{S0607}。
  \item 证书与学籍侧：学校应当严格按照招生时确定的办学类型、学习形式以及学生招生录取时填报的个人信息填写、颁发学历证书、学位证书及其他学业证书；在校期间变更姓名、出生日期等证书需填写的个人信息，应当有合理、充分的理由并提供有法定效力的相应证明文件 \nref{S0040}。
\end{itemize}
\nref{S0523}
\nref{S0182}
\nref{S0067}
\nref{S0693}
\nref{S0295}
\nref{S0607}
\nref{S0040}
\end{qa}

\iin{教育部学生服务与素质发展中心（学信网/研招网）}
\ith{系统操作}
\ist{信息与工具}

\begin{qa}{S0422}{政策网页后来改版或撤下，怎样保留能复核的原始依据？}
看到当年正式通知后，保存发布机关、页面标题、发布日期、适用届别、访问日期和原始网址；同时下载官方附件原件，保留文件名、落款页和必要的文件哈希。关键条款可记录页码/段落并保留页面打印件或截图，但截图应带网址和日期，不能只存群聊转发。之后如网页更新，重新下载并对比版本，记录变更时间；申请或申诉时引用最新有效文件，并把旧版作为变更轨迹而非当前规则。系统材料由高校提供，遇到版本冲突应找发布单位核实；涉及个人信息的截图不要公开传播。
\nref{S0761}
\qscope{操作原则适用于一般情形；具体流程、期限和便利措施以招生单位当年正式通知为准。}
\end{qa}

\section{21.10 规则版本核验}
\begin{qa}{U202610212101}{怎样快速比较两版推免细则，判断改动会不会影响自己？}
不要只看文件标题或发布日期。先从学校官网或推免服务系统下载两版原文并记录发布主体、年份、适用对象、生效时间和附件；再围绕资格条件、课程范围、综合成绩公式、名额基数、公示异议五项逐条对照。把每处改动标为“文字调整/口径变化/新增材料/时间变化”，说明对自己的影响和需向学校确认的问题。网页后续改版时保留下载原件、原始链接、访问日期及必要的校验值，旧版只能用于追踪变化，当前适用规则仍要由学校说明。\nref{S0761} \nref{S0762}
\end{qa}
\begin{qa}{U202610212102}{系统个人信息和学籍记录对不上，应该先改系统还是先提交材料？}
先区分问题字段和数据归属：姓名、证件号或学籍状态应先向本科院校学籍管理部门核实原始学籍数据；推免资格或推荐名单错误联系本校推免主管部门；目标院校报名信息则按招生单位通知联系研招办。不要在不同系统反复改写造成多套记录。整理字段名称、系统显示值、学校正式记录、页面截图和错误发生时间，分别询问负责部门由谁更正、是否需要证明材料、更新需要多久；若临近截止，先通过官方渠道提交书面说明并留回执，再按指引补齐更正结果。系统内资格审核、报名、录取及备案均有正式流程，个人说明不能替代系统记录。\nref{S0761}
\end{qa}
